%
% ---------------------------------------------------------------
% Copyright (C) 2012-2018 Gang Li
% ---------------------------------------------------------------
%
% This work is the default powerdot-tuliplab style test file and may be
% distributed and/or modified under the conditions of the LaTeX Project Public
% License, either version 1.3 of this license or (at your option) any later
% version. The latest version of this license is in
% http://www.latex-project.org/lppl.txt and version 1.3 or later is part of all
% distributions of LaTeX version 2003/12/01 or later.
%
% This work has the LPPL maintenance status "maintained".
%
% This Current Maintainer of this work is Gang Li.
%
%

\documentclass[
 size=14pt,
 paper=smartboard,  %a4paper, smartboard, screen
 mode=present, 		%present, handout, print
 display=slides, 	% slidesnotes, notes, slides
 style=tuliplab,  	% TULIP Lab style
 pauseslide,
 fleqn,leqno]{powerdot}


\usepackage{amssymb}
\usepackage{amsmath}
\usepackage{rotating}
\usepackage{graphicx}
\usepackage{boxedminipage}
\usepackage{media9}
\usepackage{rotate}
\usepackage{calc}
\usepackage[absolute]{textpos}
\usepackage{psfrag,overpic}
\usepackage{fouriernc}
\usepackage{pstricks,pst-node,pst-text,pst-3d,pst-grad}
\usepackage{moreverb,epsfig,subfigure}
\usepackage{pstricks}
\usepackage{pstricks-add}
\usepackage{pst-text}
\usepackage{pst-node, pst-tree}
\usepackage{booktabs}
\usepackage{etex}
\usepackage{breqn}
\usepackage{multirow}
\usepackage{gitinfo2}
\usepackage{xcolor}

\usepackage{todonotes}
% \usepackage{pst-rel-points}
\usepackage{animate}
\usepackage{fontawesome}

\usepackage{textcomp}

\usepackage{listings}
\lstset{frameround=fttt,
frame=trBL,
stringstyle=\ttfamily,
backgroundcolor=\color{yellow!20},
basicstyle=\footnotesize\ttfamily}
\lstnewenvironment{code}{
\lstset{frame=single,escapeinside=`',
backgroundcolor=\color{yellow!20},
basicstyle=\footnotesize\ttfamily}
}{}

%%tikz
\usepackage{verbatim}
\usepackage{amssymb}
\usepackage{amsmath}
\usepackage{tikz}
\usetikzlibrary{arrows.meta}
\tikzset{%
  >={Latex[width=2mm,length=2mm]},
  % Specifications for style of nodes:
            base/.style = {rectangle, rounded corners, draw=black,
                           minimum width=9cm, minimum height=1cm,
                           text centered, font=\sffamily},
            startstop/.style = {base, fill=red!30},
            activityStarts/.style = {base, fill=blue!30},
            process/.style = {base, minimum width=8cm, fill=orange!15,
                           font=\ttfamily},
}

\usepackage{hyperref}
\hypersetup{ % TODO: PDF meta Data
  pdftitle={Presentation Skill},
  pdfauthor={Xin Han},
  pdfpagemode={FullScreen},
  pdfborder={0 0 0}
}

%smartdiagram
\usepackage{smartdiagram}
\usetikzlibrary{shapes.geometric}


% \usepackage{auto-pst-pdf}
% package to show source code

\definecolor{LightGray}{rgb}{0.9,0.9,0.9}
\newlength{\pixel}\setlength\pixel{0.000714285714\slidewidth}
\setlength{\TPHorizModule}{\slidewidth}
\setlength{\TPVertModule}{\slideheight}
\newcommand\highlight[1]{\fbox{#1}}
\newcommand\icite[1]{{\footnotesize [#1]}}

\newcommand\twotonebox[2]{\fcolorbox{pdcolor2}{pdcolor2}
{#1\vphantom{#2}}\fcolorbox{pdcolor2}{white}{#2\vphantom{#1}}}
\newcommand\twotoneboxo[2]{\fcolorbox{pdcolor2}{pdcolor2}
{#1}\fcolorbox{pdcolor2}{white}{#2}}
\newcommand\vpspace[1]{\vphantom{\vspace{#1}}}
\newcommand\hpspace[1]{\hphantom{\hspace{#1}}}
\newcommand\COMMENT[1]{}

\newcommand\placepos[3]{\hbox to\z@{\kern#1
        \raisebox{-#2}[\z@][\z@]{#3}\hss}\ignorespaces}

\renewcommand{\baselinestretch}{1.2}


\newcommand{\draftnote}[3]{
	\todo[author=#2,color=#1!30,size=\footnotesize]{\textsf{#3}}}
% TODO: add yourself here:
%
\newcommand{\gangli}[1]{\draftnote{blue}{GLi:}{#1}}
\newcommand{\xinhan}[1]{\draftnote{green}{XH:}{#1}}
\newcommand{\gliMarker}
	{\todo[author=GLi,size=\tiny,inline,color=blue!40]
	{Gang Li has worked up to here.}}
\newcommand{\snMarker}
	{\todo[author=XH,size=\tiny,inline,color=green!40]
	{Xin Han has worked up to here.}}

%%%%%%%%%%%%%%%%%%%%%%%%%%%%%%%%%%%%%%%%%%%%%%%%%%%%%%%%%%%%%%%%%%%%%%%%
% title
% TODO: Customize to your Own Title, Name, Address
%
\title{Presentation Skill}
\author{Xin Han
\\
\\Xi'an Shiyou University
}
\date{\gitCommitterDate}


% Customize the setting of slides
\pdsetup{
% TODO: Customize the left footer, and right footer
rf=\href{http://www.tulip.org.au}{
Last Changed by: \textsc{\gitCommitterName}\ \gitVtagn-\gitAbbrevHash\ (\gitAuthorDate)
},
cf={Presentation Skill},
}


\begin{document}

\maketitle

%\begin{slide}{Overview}
%\tableofcontents[content=sections]
%\end{slide}


%%==========================================================================================
%%
\begin{slide}[toc=,bm=]{Overview}
\tableofcontents[content=currentsection,type=1]
\end{slide}
%%
%%==========================================================================================


\section{TULIP Lab \LaTeX\ Template}

\begin{slide}{powerdot-tuliplab}
\selectcolormodel{rgb}
\textit{powerdot-tuliplab} is the \LaTeX\ package used in TULIP Lab for presentation drafting.
Publishing in \textcolor{orange}{\faGithub:}\url{https://github.com/tulip-lab}.
\begin{center}
\smartdiagramset{description title width=2.2cm,
description text width=11cm,}
\smartdiagram[descriptive diagram]{
  {\textit{powerdot-tuliplab.sty},{TULIP-Lab Style File}},
  {\textit{tuliplab-P00.tex}, {A Simple Presentation Sample with Transitions}},
  {\textit{tuliplab-P01.tex},One Academic Presentation Sample with References List},
  {\textit{tuliplab-P03.tex},One Sample file with PSTrick figures.},
  {\textit{logos/tulip},Folder of images for slides graphics.}}
\end{center}
\end{slide}

\begin{slide}[toc=,bm=]{Installation \& Compiling}
\selectcolormodel{rgb}
\begin{description}
  \item[Installation] Typically this package should be installed automatically with a full \LaTeX\ package,
      or you can install directly from \textit{ctan.org}.
      \begin{itemize}
        \item Copy files the whole folder powerdot-tuliplab to your home directory under \textit{~/texmf/tex/latex/powerdot-tuliplab}
        \item Refresh \LaTeX\ platform
      \end{itemize}
  \item[Compiling]You should use the follow chain operators:\\ \begin{center}
\smartdiagram[sequence diagram]{\LaTeX\ $\to$ DVI,DVI $\to$ PS,PS $\to$ PDF}
\end{center}
\end{description}
\end{slide}
%%
%%==========================================================================================

%%==========================================================================================
%%
\begin{slide}{TempLaTeX}
\selectcolormodel{rgb}
TempLaTeX is the \LaTeX\ package used
in TULIP Lab for paper drafting.
It is published in  \textcolor{orange}{\faGithub:}
\url{https://github.com/tulip-lab/templatex/tree/master/templatex}

\begin{center}
\smartdiagramset{description title width=2.2cm,
description text width=13cm,}
\scalebox{0.7}{\smartdiagram[descriptive diagram]{
  {\textit{report.tex},{the complete report stem file}},
  {\textit{report-htm.tex}, {the complete report stem file for tourism and hospitality paper}},
  {\textit{report-acm.tex},the stem file for ACM conference proceedings},
  {\textit{report-ieee.tex},the stem file for IEEE conference proceedings},
  {\textit{report-lncs.tex},the stem file for Springer LNCS conference proceedings;},
  {\textit{slides.tex},one sample file for slides with notes},
  {\textit{poster.tex},one sample poster}}}
\end{center}
\end{slide}
%%
%%==========================================================================================

%%==========================================================================================
%%
\begin{slide}[toc=,bm=]{Set \& Compiling}
You are expected to use git to host the repository,
and use \textit{git-flow} to manage
the collaborative writing platform
for academic paper authoring.

\begin{center}
\selectcolormodel{rgb}
\smartdiagramset{module shape=diamond,
font=\scriptsize,
module minimum width=1cm,
module minimum height=1cm,
text width=1cm,
circular distance=2cm,
circular final arrow disabled=true,
}
\smartdiagram[circular diagram:clockwise]{git clone,Update self info,Set gitinfo2,Compile needed}
\end{center}
\end{slide}
%%
%%==========================================================================================

%%==========================================================================================
%%
\begin{slide}[method=direct]{Update Self Information}
Open the \textit{preamble.tex},
and update with your own information at the following places:

\begin{itemize}
\item
Add yourself:
\begin{lstlisting}[language=TeX]
\newcommand{\yourname}[1]{\draftnote{blue}{[You: #1]}}
% other colors include blue, red, purple, cyan, darkgreen, etc.
\newcommand{\gliMarker}
	{\todo[author=GLi,size=\tiny,inline,color=blue!40]
	{Gang Li has worked up to here.}}
\end{lstlisting}

\item
Fill the proper information:
\begin{lstlisting}[language=TeX]
\hypersetup
{
    pdfauthor={\gitAuthorName},
    pdfsubject={TULIP Lab},
    pdftitle={},
    pdfkeywords={TULIP Lab, Data Science},
    bookmarks=true,
}
\end{lstlisting}
\end{itemize}
\end{slide}
%%
%%==========================================================================================

%%==========================================================================================
%%
\begin{slide}[toc=,bm=,method=direct]{Update Self Information}
Open the \textit{preamble.tex},
and update with your own information at the following places:
\begin{itemize}
  \item Customize Title, Name, Address:
  \begin{lstlisting}[language=TeX]
  \title{Presentation Title}
    \author{Gang Li
    \\
    Deakin University
    % \href{mailto:gangli@acm.org}{gangli@acm.org}
    % \and % more authors
    }
    \end{lstlisting}
  \item Customize the setting of slides:
  \begin{lstlisting}[language=TeX]
  rf=\href{http://www.tulip.org.au}{
  Last Changed by: \textsc{\gitCommitterName}\
  \gitVtagn-\gitAbbrevHash\ (\gitAuthorDate)},
  cf={Presentation Title},}
  \end{lstlisting}
\end{itemize}
\end{slide}
%%
%%==========================================================================================

%%==========================================================================================
%%
\begin{slide}[method=direct]{Set Gitinfo2}
The package contains a gitinfo folder
with three executive bash shell scripts:
\textcolor{orange}{post-checkout},
\textcolor{orange}{post-commit} and
\textcolor{orange}{post-merge}.

\begin{itemize}
    \item  Copy or move these three files
           into your repository's \textit{.git/hooks} folder.
    \item Test Run: check out or pull your repository,
          and it should generate/update the file \textit{./git/gitHeadInfo.gin}
          in your local repository.
    \item Compile the main \LaTeX\ file, and view the PDF.
\end{itemize}

\end{slide}
%%
%%==========================================================================================

%%==========================================================================================
%%
\begin{slide}[method=direct]{Basic Usage}
\begin{itemize}
  \item Document setting:
  \begin{lstlisting}[language=TeX]
  \documentclass[
     size=14pt,
     paper=smartboard,  %a4paper, smartboard, screen
     mode=present, 		%present, handout, print
     display=notes, 	% slidesnotes, notes, slides
     style=tuliplab,  	% TULIP Lab style
     pauseslide,
     fleqn,leqno]{powerdot}
 \end{lstlisting}
 The \textit{print mode} can be used to print the slides,
 it deletes the overlays and transition effects.
 Other modes are \textit{present},
 which is the default mode for presentations;
 and \textit{handout} which
 produces a black and white overview of the slides,
 printing two slides per page.
 \end{itemize}
 \end{slide}
%%
%%==========================================================================================

%%==========================================================================================
%%
\begin{slide}[toc=,bm=,method=direct]{Basic Usage}
\begin{itemize}
  \item Section \& Slide:
  \begin{lstlisting}[language=TeX]
    % section: title takes up full slide
    \section{First section}
    \begin{slide}{Slide Title}
        \begin{itemize}
        \item This is an item
        \item Second item
        \item Third item
        \end{itemize}
    \end{slide}
  \end{lstlisting}
  Pretty much anything you can do in a normal \LaTeX\
      document:figures,
      tables,
      equations,
      normal text~{et.~al.}
      between \textit{\textbackslash begin\{slide\}
      and \textbackslash end\{slide\}}
  \end{itemize}
  \end{slide}
%%
%%==========================================================================================

%%==========================================================================================
%%
\begin{slide}[toc=,bm=,method=direct]{Basic Usage}
 \begin{itemize}
  \item Adding notes:
  \begin{lstlisting}[language=TeX]
  \documentclass[
     size=14pt,
     paper=smartboard,  %a4paper, smartboard, screen
     mode=present, 		%present, handout, print
     display=notes, 	% slidesnotes, notes, slides
     style=tuliplab,  	% TULIP Lab style
     pauseslide,
     fleqn,leqno]{powerdot}
    %
    \begin{slide}{Slide Title}
        \begin{itemize}
        \item This is an item
        \end{itemize}
    \end{slide}
    \begin{note}{About items}
        Mention that lists of items can be customised.
    \end{note}
  \end{lstlisting}
  \textit{display = slidesnotes} print the notes and the slides,
  while \textit{display = slides} prints only the slides.
\end{itemize}
\end{slide}
%%
%%==========================================================================================
\begin{slide}{TULIP Presentation Procedure}
\selectcolormodel{rgb}
\begin{itemize}
  \item General presentation procedure in TULIP:
\begin{center}

\scalebox{0.6}{
\begin{tikzpicture}[node distance=2cm,
    every node/.style={fill=white, font=\sffamily}, align=center]
  \node (start)             [activityStarts]          {Activity starts};
  \node (groupMeeting)      [process, below of=start] {Report process in $\infty$ meeting};
  \node (preparProblem)     [process, below of=groupMeeting]{Prepar atleast 10 problems};
  \node (readingTeam)       [process, below of=preparProblem]{Report in Reading Team};
  \node (collectProblem)    [process, below of=readingTeam]{Collect Problems $\&$ Answser};
  \node (tulipSeminar)      [process, below of=collectProblem]{Report in Tulip Seminar};
  \node (collectProblem2)   [process, below of=tulipSeminar]{Collect Problems $\&$ Answser};
  \node (onResumeBlock)     [startstop, below of=collectProblem2]  {Report in External Place};

  \draw[->]      (start) -- node{Use PPT}(groupMeeting);
  \draw[->]     (groupMeeting) -- (preparProblem);
  \draw[->]     (preparProblem) -- node{Use Slides}(readingTeam);
  \draw[->]     (readingTeam) -- (collectProblem);
  \draw[->]     (collectProblem) -- (tulipSeminar);
  \draw[->]     (tulipSeminar) -- (collectProblem2);
  \draw[->]     (collectProblem2) -- (onResumeBlock);

  \end{tikzpicture}}
\end{center}
\end{itemize}

\end{slide}
%%==========================================================================================
%%

%%
%%==========================================================================================

\section{What \& How to Present}

%%==========================================================================================
%%
\begin{slide}[toc=,bm=]{What and How to Present}
Choice and organization of the material to be presented.
\begin{itemize}
  \item \textcolor{orange}{Communicate the Key Ideas}\\
  Make sure that your talk emphasizes the key ideas and
  skips over what is standard,
  obvious,
  or merely complicated.
  \item \textcolor{orange}{Don't get Bogged Down in Details}\\
  Details are out of place in an oral presentation.
  This rule cannot be over-emphasized
  \item \textcolor{orange}{Structure Your Talk}\\
  Your presentation should be broken into several distinct parts,
  each with its own objectives and style.

\end{itemize}
\end{slide}
%%
%%==========================================================================================

%%==========================================================================================
%%
\begin{slide}[toc=,bm=]{What and How to Present}
Choice and organization of the material to be presented.
\begin{itemize}

  \item \textcolor{orange}{Know Your Audience}\\
  Think through the average level of expertise
  in your audience and present your results accordingly
  \item \textcolor{orange}{Use an Organized Approach}\\
  A simple three-part template for producing a talk is as follows:
  \begin{itemize}
    \item Introduction
    \item Body
    \item Conclusion
  \end{itemize}
  \end{itemize}
\end{slide}
%%
%%==========================================================================================

%%==========================================================================================
%%
\begin{slide}{The Introduction}
It sets the tone for the entire talk.
It determines whether the audience will prick up their ears,
or remain slumped in their chairs.
First impressions are very important.
\begin{itemize}
  \item \textcolor{orange}{Define the Problem}\\
  No matter how difficult and technical the problem;
  it can usually be described succinctly and
  accurately in less than five minutes.
  \item \textcolor{orange}{Motivate the Audience}\\
  Explain why the problem is so important.
  How does the problem fit into the larger picture?
  What are its applications?
  What makes the problem nontrivial?
  \item \textcolor{orange}{Introduce Terminology}\\
  All terms must be introduced early.
\end{itemize}
\end{slide}
%%
%%==========================================================================================

%%==========================================================================================
%%
\begin{slide}[toc=,bm=]{The Introduction}
\begin{itemize}

    \item \textcolor{orange}{Discuss Earlier Work}\\
      Compared and explain why your work is different
      from past work.

    \item \textcolor{orange}{Emphasize the Contributions of your Paper}\\
      Make sure that you explicitly and succinctly state the
      contributions made by your paper.
      Often it is the only thing that audience carry
      away from the talk.

    \item \textcolor{orange}{Consider putting your Conclusion in the Introduction}  \\
        Take a bold step and put a short conclusion
        in the introduction.
        Let everyone know up front the direction
        you are headed so that the audience can
         focus on the details.
\end{itemize}
\end{slide}
%%
%%==========================================================================================

%%==========================================================================================
%%
\begin{slide}{The Main Body}
This contains the meat of your presentation,
and is the point at which the attention of
the audience will start to waver
if you messed up your Introduction.
Consider the following suggestions.
\begin{itemize}
  \item \textcolor{orange}{Abstract the Major Results.}\\
  Describe the key results of the presentation.
  You will probably have to get a little technical
  here but do so gradually and carefully.

  \item  \textcolor{orange}{Explain the Significance of the Results}\\
  Focus on anything unexpected or crucial to supporting your conclusions.

  \item  \textcolor{orange}{Sketch a Proof of the Crucial Results}\\
   State the hypotheses and experimental design
   as simply as possible.
\end{itemize}
\end{slide}
%%
%%==========================================================================================

%%==========================================================================================
%%
\begin{slide}[toc=,bm=]{The Body}
\begin{itemize}
  \item \textcolor{orange}{Use Props and Pictures}\\
  ``Things seen are mightier than things heard.''
  Be sure to add figures and photos to your slides,
  where appropriate.

  \item \textcolor{orange}{Don't Inflict Pain on the Audience}\\
  Never waste an audience's time taking them
  through a step-by-step history of your project.

  \item \textcolor{orange}{Avoid Complex Tables}\\
  Don't cram a lot of numbers
  into small tables on your slides. Never.

  \item \textcolor{orange}{Have a Purpose and Conclusion}\\
  Consider putting a written conclusion at
  the bottom of each key slide.
  Each slide should have a point that is being made.
  \end{itemize}
\end{slide}
%%
%%==========================================================================================

%%==========================================================================================
%%
\begin{slide}{The Conclusion}
Your aim here is to restate the lessons learned in a short,
concise manner.
\begin{itemize}
  \item \textcolor{orange}{Hindsight is Clearer} than Foresight\\
  You can now make observations that
  would have been confusing if they were introduced earlier.
  Use this opportunity to refer to statements that
  you have made in the previous three sections and
  weave them into a coherent synopsis.

  \item \textcolor{orange}{Indicate that your Talk is Over}\\
  An acceptable way to do this is to say ``Thank you.
  Are there any questions?''

\end{itemize}
\end{slide}

\section{Question Time}

\begin{slide}[toc=,bm=]{Question Time}
\begin{description}
  \item[Genuine request] It should cause you no difficulty if you are adequately prepared.\pause
  \item[Selfish question] It is \textcolor{orange}{polite to take a few seconds} to compose a reply that directly or indirectly compliments the questioner.\pause
  \item[Malicious question]Be \textcolor{orange}{prepared},
  be polite,
  and avoid getting involved in a \textcolor{orange}{lengthy exchange}.
\end{description}
\end{slide}

\section{Getting Through to the Audience}
\begin{slide}[toc=,bm=]{Getting Through to the Audience}
The next major hurdle is when you find yourself actually
standing in front of the audience.
Faulty delivery can ruin even a wellprepared talk.

\begin{itemize}
  \item \textcolor{orange}{Practice Your Talk}\\
   The practice must be verbal, not just mental.
   For many reasons,
   a couple verbal practice rounds is one of the most
   effective ways to improve your communication effectiveness.

  \item \textcolor{orange}{Use Repetition}\\
  The tried and true strategy for presentations is to:
  ``Tell them what you're going to tell them (the Introduction).
  Tell them (the Body).
  And then tell them what you told them (the Conclusion).''

  \item \textcolor{orange}{Convey Enthusiasm, Excitement,
  Confidence}\\
  Believe in what you are doing.
  Let the audience know that you have something very
  interesting to share with them.
\end{itemize}
\end{slide}

\begin{slide}[toc=,bm=]{Getting Through to the Audience}
\begin{itemize}
\item \textcolor{orange}{Control Your Voice}\\
  Speak clearly and with sufficient volume.
  Don't speak in a monotone.
  Avoid utterances.
  Avoid fashionable turns of phrase.
  Avoid hype.

  \item \textcolor{orange}{Minimize Language Difficulties}\\
  It is a good idea to
  get a native speaker to look over your
  transparencies before you deliver the talk.
  Use a prepared text,
  if necessary.

  \item \textcolor{orange}{Don't Start Your Talk with an Apology}\\
  The speaker may be doing nothing more than
  trying to deal with his/her own nervousness,
  but as a result the audience now has a lower
  opinion and expectation of the speaker and their presentation.
\end{itemize}
\end{slide}

\section{Visual and Aural Aids}
\begin{slide}[toc=,bm=]{Visual and Aural Aids}
Now that you have a well-prepared talk and can deliver it with style.
\begin{itemize}
  \item \textcolor{orange}{Use Color Effectively}\\
  The use of color can enhance a presentation,
  particularly when used in moderation in the
  figures and diagrams.
  Only use high contrast color combinations.

  \item \textcolor{orange}{Use Pictures and Tables}\\
  Use photos and drawings whenever possible,
  but only when they illustrate the point
  you are trying to make.

  \item \textcolor{orange}{Familiarize Yourself with the Stage} \\
  Introduce yourself to the session chair if you haven't already.
  If you have 35mim slides,
  be sure they are mounted properly and are in the correct order.
\end{itemize}
\end{slide}
%%==========================================================================================
%
\begin{slide}[toc=,bm=]{Questions?}
\begin{center}
\begin{figure}
    \animategraphics[autoplay, loop, height=0.5\textheight]{5}{figures//management//gif//question//q_}{1}{30}
\end{figure}
\end{center}
\end{slide}
%%
%%==========================================================================================


%%==========================================================================================
% TODO: Contact Page
\begin{wideslide}[toc=,bm=]{Contact Information}
\centering
\vspace{\stretch{1}}
\twocolumn[
lcolwidth=0.35\linewidth,
rcolwidth=0.65\linewidth
]
{
% \centerline{\includegraphics[scale=.2]{tulip-logo.eps}}
}
{
\vspace{\stretch{1}}
Xin Han\\
School of Computer Science\\
Xi'an Shiyou University, China

\begin{description}
 \item[\textcolor{orange}{\faEnvelope}] \href{mailto:xin.han@tulip.org.au}
 {\textsc{\footnotesize{xin.han@tulip.org.au}}}

 \item[\textcolor{orange}{\faHome}] \href{http://www.tulip.org.au}
 {\textsc{\footnotesize{Team for Universal Learning and Intelligent Processing}}}
\end{description}
}
\vspace{\stretch{1}}
\end{wideslide}

\end{document}

\endinput
